\chapter{Рантайм и инженерия}
\label{ch:runtime}

До сих пор мы описывали \emph{что} делают операторы. Эта глава — про то, \emph{как}
они реально запускаются, приостанавливаются, где хранятся данные и во что это
обходится по памяти и времени. Это инженерный слой, по которому предстоит работать,
поэтому здесь мы конкретны, а оценки даём даже там, где они лишь порядковые.

\section{Планировщик как конечный автомат}
\label{sec:scheduler}

Никакого «мистического решающего разума», дёргающего операторы, нет. Есть
\emph{цикл событий} с приоритетным планировщиком. Управление двухуровневое
(глава~\ref{ch:walker}): грубый \emph{режим} задаётся конечным автоматом, а внутри
режима \emph{диспетчер} раздаёт операции по бюджету $\rho$.

\begin{definition}[Единственный дефицитный ресурс]
Почти всё в системе дёшево (обход, статистика, распад, векторный поиск — миллисекунды,
раздел~\ref{sec:costs}). \emph{Дорог} только один ресурс — \textbf{слот
LLM-инференса} (один GPU, вызовы сериализуются). Поэтому планировщик по сути
распределяет \emph{слоты LLM-вызовов} между передним планом (ответ) и фоном
(fusion, метакогниция, переваривание). Всё остальное диспетчеру раздать легко.
\end{definition}

\subsection{Режимы (состояния автомата)}

\begin{center}
\begin{tabular}{p{0.16\textwidth} p{0.10\textwidth} p{0.60\textwidth}}
\toprule
\textbf{Режим} & $\rho$ & \textbf{Что делает диспетчер} \\
\midrule
\textsc{Engaged} & высок & приоритет — retrieval $+$ композиция ответа; фон приглушён/приостановлен \\
\textsc{Idle} & $\approx 0$ & фоновый разведчик (любопытство $+$ лёгкий fusion), лёгкая консолидация \\
\textsc{Sleeping} & $\approx 0$, долго & телепорт-скраб, тяжёлая консолидация, репетиция важного \\
\textsc{Reflecting} & сред. & объяснение сработавшей аномалии (метакогниция); под-режим \textsc{Idle} \\
\bottomrule
\end{tabular}
\end{center}

\subsection{Переходы}

\begin{lstlisting}[style=pseudo]
on user_message:            # экзогенное прерывание, наивысший приоритет
    freeze_background()     # заморозить текущее фоновое рассуждение (раздел 18.2)
    mode <- ENGAGED
on response_delivered and no pending_foreground:
    mode <- IDLE
on idle_timeout:            # долгий простой
    mode <- SLEEPING
on detector_fired(anomaly) and salient:
    mode <- REFLECTING      # вернётся в прежний режим по завершении
on any_state + user_message:
    preempt -> ENGAGED
\end{lstlisting}

Ключевой момент — \textbf{смещение фокуса «любопытная Рин $\to$ ответить юзеру» это
рефлекс, а не решение LLM}: сообщение пользователя — высоко-salience экзогенное
событие, оно вытесняет фон как аппаратное прерывание ($\rho$ подскочил $\Rightarrow$
фон приглушён). LLM подключается лишь на настоящих развилках суждения (бросить нить
или припарковать; вставить ассоциацию или остаться в теме) — это \emph{эскалация} к
делиберации (System~2), редкая. Рефлексы (System~1) закрывают подавляющее
большинство запусков.

\subsection{Реализация}

Простейшее — однопоточный цикл событий (события: \texttt{user\_msg},
\texttt{timer\_tick}, \texttt{detector\_fire}, \texttt{op\_complete}). Для
параллелизма — небольшой пул: поток переднего плана (ответ) и фоновый поток
(разведка), сериализуемые на общем ограничении — числе доступных LLM-слотов (так
сделано и в ReplicantLife у Toy и др.). То есть настоящий планировщик планирует
\emph{GPU-слоты}, а не «мысли».

\begin{openq}[Граница эскалации к LLM]
Какие именно решения остаются рефлексом, а какие эскалируют к LLM-делиберации, —
подвижная граница. Слишком много эскалаций $\Rightarrow$ LLM-бутылочное горло и
тормоза; слишком мало $\Rightarrow$ Рин ведёт себя механически. Настраивается на
стенде.
\end{openq}

\subsection{Произвольное внимание: внутренний forensic-запрос в Reflecting}
\label{sec:reflecting}

В таблице режимов \textsc{Reflecting} помечен как «объяснение аномалии», но это
недосказано: тут прячется механизм, которого в мануале не было явно — \emph{как Рин
думает сама с собой, без внешнего триггера}. Всё внимание до сих пор было экзогенным
(сообщение юзера как прерывание). \textsc{Reflecting} — это \emph{эндогенное,
произвольное} внимание, и оно сшивает три уже описанных куска:
\begin{enumerate}\itemsep2pt
\item \textbf{Триггер (метакогниция).} Дешёвые детекторы на кольцевом буфере телеметрии
(гл.~\ref{ch:mood}) ловят системный сбой — скачок \texttt{need\_to\_revise}, конфликт
уверенностей, затяжную руминацию.
\item \textbf{Переход (планировщик).} Автомат уходит в \textsc{Reflecting}
(\texttt{detector\_fired} выше).
\item \textbf{Механизм (forensic-ретрив).} Рин \emph{сама себе} кидает направленный
\emph{forensic}-запрос (раздел~\ref{sec:retrievaltax}) — не ассоциативный обход, а точный
подъём дерева обоснований по провенансу через query-planner — и запускает ретракцию лжи
по TMS (гл.~\ref{ch:epistemics}). Это и есть моделируемый внутренний монолог: не «что
всплыло», а «на чём это стоит и не рухнуло ли основание».
\end{enumerate}

\begin{pitfall}[Тормоза руминации обязательны]
Без ограничителей эндогенное внимание уходит в \emph{self-fulfilling-loop}: детектор
руминации сам себя кормит, Рин зависает в бесконечном самоаудите (это тот же риск, что и
самоисполняющаяся тревога, гл.~\ref{ch:mood}). Поэтому \textsc{Reflecting} жёстко
гейтируется: \emph{рейт-лимит} (не чаще, чем раз в такое-то окно), \emph{event-driven
каденция} (только на \emph{значимую} аномалию, не на каждый микроконфликт уверенностей)
и \emph{бюджет} на глубину\,/\,длительность рефлексии. Выгорел бюджет — вышли из режима,
пометив вопрос как открытый, а не крутимся дальше.
\end{pitfall}

\section{Приостановка и возобновление рассуждений}
\label{sec:suspend}

Когда сообщение пользователя вытесняет фоновое рассуждение (Рин была на середине
fusion-цепочки или любопытной прогулки), состояние надо \emph{заморозить}, чтобы потом
продолжить с того же места, а не с нуля.

\subsection{Что такое «состояние рассуждения»}

Фоновая задача — это walker (или рой) с: текущим \emph{фронтом} (множество активных
узлов $+$ их активации), \emph{затравкой/целью} (что искали), \emph{частичными
результатами} (уже собранные кандидаты, поставленные \texttt{need\_to\_discover}) и
\emph{курсором} (где остановились). Заморозка — это сериализация \emph{continuation-
токена}:

\begin{lstlisting}[style=pseudo]
Continuation = {
    task_id, goal/seed,
    frontier: [(node_id, activation), ...],   # компактно: id + float
    partial_candidates: [node_id, ...],
    pending: [need_to_discover, ...],
    cursor, priority, timestamp
}
\end{lstlisting}

Токен мал (id-шники $+$ числа) — это как захват состояния корутины при yield или
context-switch процесса в ОС (сохранить «регистры» = фронт $+$ курсор). Кладётся в
очередь приостановленных задач.

\subsection{Возобновление и устаревание}

При возврате в \textsc{Idle} задача достаётся из очереди и продолжается. Но граф мог
\emph{измениться} за время паузы (появились узлы, активации распались, что-то
спрунено), поэтому возобновление начинается с \emph{ре-валидации фронта}: выкидываем
исчезшие/распавшиеся узлы, восстанавливаем активации живых, продолжаем обход.

Не всякую задачу стоит возобновлять: пока Рин отвечала, разговор мог уйти, и нить
устарела. Поэтому у приостановленных задач своё \emph{устаревание}: если не
возобновлена за окно или вытеснена более свежей — она \emph{не удаляется бесследно}, а
роняется в граф как \emph{«недодуманная нить»} (узел с \texttt{need\_to\_discover}),
которую позже может подобрать телепорт-скраб во сне (глава~\ref{ch:sleep}). Так
прерванная мысль не теряется, а становится зацепкой для фоновой разведки.

\begin{intuition}[Замороженная мысль — это savepoint, а не труп]
Полезно думать о continuation-токене как о savepoint базы данных или о зелёном
потоке, снятом с исполнения: он хранит ровно столько, чтобы продолжить, и ничего
лишнего. А брошенная нить не исчезает — она оседает в память как «я не додумала про
$X$», и это честно по-человечески.
\end{intuition}

Каждая пара приостановка/возобновление — \emph{событие телеметрии} (глава~\ref{ch:mood}):
паттерн «меня постоянно прерывают на середине мысли» \emph{детектируется} метакогницией
как аномалия.

\section{Хранилище и индексирование}
\label{sec:storage}

Разберём предметно, из чего состоит хранилище и как оно индексируется, — и заодно
ответим, заменяет ли многоуровневое хранилище простое «окно внимания».

\subsection{Из чего состоит хранилище}

\begin{itemize}\itemsep1pt
\item \textbf{Node store} — ключ-значение \texttt{node\_id} $\to$ \{content,
embedding, метаданные\}; на диске (встраиваемое KV-хранилище: LMDB/SQLite и т.п.).
\item \textbf{Edge store} — списки смежности \texttt{node\_id} $\to$
[(target, type, weight, sign), \dots].
\item \textbf{Vector index} — индекс по эмбеддингам для поиска похожих (HNSW/FAISS),
\textbf{резидентный в RAM}; держит эмбеддинги \emph{горячего/тёплого} множества и
служит механизмом \emph{поиска якорей} по запросу (глава~\ref{ch:walker}). Важно:
\emph{поиск идёт по этому RAM-индексу, а не по диску} — на диске лежит лишь
персистентная копия эмбеддингов (backing store) и cold-тир.
\item \textbf{Структурные индексы} — по \texttt{layer}, по \texttt{importance} и т.п.,
для планировщика и распада.
\end{itemize}
Индексов, стало быть, два по сути: \emph{векторный} (семантическая близость, чтобы
найти вход) и \emph{графовая смежность} (чтобы идти от входа). Retrieval $=$ векторный
поиск якорей $\to$ обход графа от них.

\subsection{Attention-radii против тиров: это не одно и то же}

Твоя догадка про «окно/радиусы внимания» верна, но покрывает лишь половину.

\emph{Attention-radius} — это \textbf{транзиентная} граница на \emph{обход}: сколько
хопов walker расширяется от якорей за один запрос. Она держит \emph{рабочее множество}
маленьким \emph{в момент запроса}, независимо от размера графа.

\emph{Тиры} (hot/warm/cold) — это \textbf{персистентная} организация \emph{хранилища},
живущая между запросами.

Почему нужно и то, и другое: attention-radius ограничивает \emph{прогулку}, но \emph{не}
ограничивает \emph{векторный индекс якорей}. Если эмбеддинги \emph{всего} графа лежат
в одном плоском индексе, стоимость поиска якорей и RAM растут с общим числом узлов.
Тиры выносят холодные узлы \emph{из} горячего индекса (они достижимы лишь через
graph-указатели), удерживая индекс маленьким. Плюс это режет \emph{шум} ретрива: в
огромном плоском индексе косинусная близость тащит семантически близкое, но
контекстно нерелевантное (ровно жалоба Toy и др.).

\begin{center}
\begin{tabular}{p{0.22\textwidth} p{0.34\textwidth} p{0.34\textwidth}}
\toprule
 & \textbf{attention-radius} & \textbf{тиры (hot/warm/cold)} \\
\midrule
природа & транзиентная (на запрос) & персистентная (хранилище) \\
ограничивает & глубину обхода (working set) & размер индекса якорей, RAM, шум \\
когда хватает одного & если граф \emph{мал} & если граф \emph{растёт вечно} \\
\bottomrule
\end{tabular}
\end{center}

Итог: \textbf{малый граф — attention-radius достаточно, тиры не нужны; большой/растущий
— нужны оба.} Тир cold хранит огрублённый gist (глава~\ref{ch:timescales}): в warm-
индексе остаётся узел-указатель домена («знаю цепи Маркова; детали в архиве»), полная
деталь лежит в cold сжатой и \emph{подкачивается} по требованию (это и есть «чуть
дольше добывать»), как paging в ОС.

\begin{openq}[Политика тиринга]
Когда демоутить узел в cold, насколько огрублять warm-gist, когда пейджить деталь
обратно — параметры, которые надо мерить (глава~\ref{ch:experiment}). По голому диску
(раздел~\ref{sec:costs}) тиринг не обязателен; он оправдан скоростью и чистотой
ретрива, а не экономией места.
\end{openq}

\subsection{Кто к кому обращается: путь запроса и стоимости}
\label{sec:dataflow}

Ключевое уточнение, снимающее частую путаницу: \textbf{по эмбеддингам на диске никто
не ищет}. Поиск идёт по векторному индексу \emph{в RAM}; диск — это backing store и
cold-тир. Горячий путь диск почти не трогает. Ниже — кто, когда и куда обращается.

\paragraph{Путь чтения («жизнь запроса»).}
\begin{enumerate}\itemsep1pt
\item \textbf{Эмбеддинг запроса}: эмбеддер (GPU/CPU), текст $\to$ вектор. $\sim$1--10\,мс.
\item \textbf{Поиск якорей}: вектор $\to$ индекс HNSW \emph{в RAM} $\to$ top-$k$
\texttt{node\_id}. $\sim$мс (sub-10\,мс даже при $10^6$). \emph{Диска нет.}
\item \textbf{Обход}: от якорей по igraph \emph{в RAM} (ограниченный радиус,
mood-bias) $\to$ активированная окрестность. $\sim$мс--10\,мс. \emph{Диска нет.}
\item \textbf{Достать контент}: для горстки финальных узлов (5--50) текст из
node-store — RAM-кэш (горячее) или SSD (несколько KV-чтений, sub-мс, обычно уже в OS
page cache). Пренебрежимо.
\item \textbf{Композиция}: контекст $\to$ LLM. GPU-VRAM. $\sim$1--3\,с.
\textbf{Доминирует.}
\end{enumerate}
То есть всё, кроме LLM, — в RAM и $\sim$мс; диск трогается лишь чтобы достать
\emph{текст} нескольких финальных узлов (дёшево), а \emph{не} для поиска.

\paragraph{Путь записи (новый узел / fusion).} Контент+эмбеддинг+метаданные $\to$
node-store (диск, асинхронно); эмбеддинг $\to$ insert в RAM-индекс (HNSW $\sim$мс,
амортизированно); узел/рёбра $\to$ в igraph (мутации батчить, глава~\ref{ch:experiment}).
Дёшево.

\paragraph{Cold-путь (единственная «дисковая» латентность, и она редкая).} Запросу
нужна архивная деталь $\to$ в warm-индексе есть узел-указатель домена $\to$ вторая
ступень: чтение из cold-store (SSD) $+$ декомпрессия $\to$ подкачка в рабочую память.
$\sim$мс--десятки\,мс. Редко, потому что hot/warm держим в RAM.

\paragraph{Старт.} При загрузке структура графа и векторный индекс поднимаются или
mmap-ятся в RAM — одноразово, не на каждый запрос.

\begin{center}
\begin{tabular}{p{0.05\textwidth} p{0.30\textwidth} p{0.22\textwidth} p{0.28\textwidth}}
\toprule
\# & \textbf{Операция} & \textbf{Ярус} & \textbf{Стоимость} \\
\midrule
1 & эмбеддинг запроса & GPU/CPU & $\sim$1--10\,мс \\
2 & поиск якорей (HNSW) & \textbf{RAM} & $\sim$мс \\
3 & обход (igraph) & \textbf{RAM} & $\sim$мс--10\,мс \\
4 & контент финальных узлов & RAM-кэш / SSD & sub-мс, пренебр. \\
5 & композиция (LLM) & \textbf{GPU-VRAM} & \textbf{$\sim$1--3\,с (узкое место)} \\
— & подкачка cold-детали & SSD & $\sim$мс--десятки\,мс (редко) \\
\bottomrule
\end{tabular}
\end{center}

\section{Event sourcing: неизменяемый лог под изменяемым графом}
\label{sec:eventlog}

Граф памяти \emph{не должен} быть единственным хранилищем прошлого. Лучше три уровня:
\begin{itemize}\itemsep1pt
\item \textbf{Event log} — неизменяемая, только-на-дозапись история всего, что
случилось: наблюдения, коммиты, ретракции, исправления, удаления, консолидации.
\item \textbf{Memory graph} — текущая материализованная модель памяти,
\emph{перестраиваемая проекция} поверх лога.
\item \textbf{Indexes} — производные структуры для поиска (векторный, forensic).
\end{itemize}
Запись лога: \{timestamp, тип (observed/committed/retracted/consolidated/deleted/
corrected), payload, source\_context, owner\}. Важно: \texttt{source\_context} и
\texttt{owner} \emph{рождаются здесь}, в момент захвата — отсюда и берётся провенанс
(глава~\ref{ch:substrate}).

Зачем это несущее: debug («почему Рин так решила» — реплеем лога), rollback,
пересборка графа после смены схемы, \textbf{сравнение версий операторов} — поменял
алгоритм консолидации, пересобрал проекцию из \emph{того же} лога, сравнил (идеально
ложится на стенд, глава~\ref{ch:experiment}), удаление, миграции.

\begin{openq}[Рост лога и компакция]
Неизменяемый лог растёт неограниченно. На старте — хранить и не думать (разрастётся —
решим). Дальше нужна политика: периодический \emph{снапшот} проекции $+$ обрезка
старого лога, либо тиринг лога в cold (раздел~\ref{sec:storage}). Точный порог и
формат снапшота — открыты.
\end{openq}

\subsection{Снапшот-чтения: почему обходчик и коммиттер не сталкиваются}
\label{sec:snapshotreads}

Раз обходчиков много и они работают параллельно с коммиттером, напрашивается страх
гонок: RW читает узел, который коммиттер как раз меняет — не получит ли он рваное,
полу-обновлённое состояние? При нашей раскладке ролей (walker читает и предлагает,
структурно пишет только сериализованный коммиттер, гл.~\ref{ch:walker},
раздел~\ref{sec:proposecommit}) страх почти растворяется, и event sourcing закрывает
остаток.

Ключ в том, что «изменение» узла у нас — \emph{не} правка на месте, а \emph{новая
версия}: коммиттер дописывает событие в лог и обновляет проекцию, старая версия узла
никуда не пропадает мгновенно. Значит обходчики читают \emph{консистентный снапшот}
проекции (MVCC-подобно): RW2, зашедший в «меняющийся» узел, видит его до-коммитную
версию целиком — рваных чтений нет, потому что нет правки на месте, которую можно
застать на середине. Read $\leftrightarrow$ read не конфликтует никогда; read
$\leftrightarrow$ write решается версионностью; write $\leftrightarrow$ write не
существует по построению (писатель один). Остаётся лишь мягкое побочное состояние —
рабочая яркость $a_{\text{work}}$ и телеметрия внимания; оно \emph{коммутативно} (для
яркости — через максимум, не сумму, гл.~\ref{ch:walker}), поэтому порядок конкурентных
подсветок неважен и блокировки не нужны.

\begin{intuition}[Гонку не арбитрируют — её конструктивно исключают]
Мы не пишем сложный протокол блокировок, кто кого пропускает. Мы устроили систему так,
что сталкиваться \emph{нечему}: один структурный писатель, версионные чтения для
всех остальных, коммутативная телеметрия. Это дешевле и надёжнее любого арбитража
постфактум — та же философия, что «противоречие лучше не разрешать задним числом, а
не допускать молчаливой перезаписи» (гл.~\ref{ch:sleep}).
\end{intuition}

\subsection{Устаревшие предложения: коммиттер валидирует, а не просто пишет}
\label{sec:staleproposals}

Снапшот-чтения убирают гонку на \emph{чтении}, но остаётся тонкий изъян на
\emph{записи}. Обходчик мог посчитать предложение на версии графа
\texttt{snapshot=123}, а к моменту коммита проекция уже на \texttt{129}, где часть
посылок успела устареть: узел-родитель ретрактнут, слит, понижен в уверенности,
изменил \texttt{layer} или \texttt{visibility}. Если коммиттер применит предложение
вслепую, гонку по данным мы убрали, но получили \emph{семантически устаревший коммит} —
что не лучше.

Поэтому предложение — не голая операция, а операция с \emph{контекстом валидности}:
\{\texttt{op}, \texttt{members}, \texttt{read\_snapshot\_id}, \texttt{parents},
\texttt{proposal\_ttl}, \texttt{validity\_conditions}\}. Коммиттер перед применением
\emph{валидирует}, а не сразу пишет: живы ли ещё родители? не deprecated? не изменились
ли их \texttt{confidence}\,/\,\texttt{layer}\,/\,\texttt{visibility} с момента
\texttt{read\_snapshot\_id}? не вышел ли \texttt{proposal\_ttl}? По итогу — одно из
четырёх: \emph{применить} (условия держатся), \emph{rebase} (пересчитать на текущей
версии, если посылки лишь сдвинулись), \emph{downgrade} (принять как более слабую
гипотезу) или \emph{reject} (посылки развалились). Это оптимистический контроль
конкуренции: не блокируем заранее, а проверяем на коммите. Без этого read/write-гонку
уберёшь технически, а память тихо наполнится решениями, принятыми по уже неверным
предпосылкам.

\subsection{«Забудь это»: удаление при неизменяемом логе}
\label{sec:forget}

Неизменяемый лог порождает очевидный конфликт: а что происходит, когда пользователь
говорит «Рин, забудь последние 10 минут»? Для локального компаньона это не юридическая
занудность, а конкретный будущий UX, и он должен что-то значить.

Ключ к разрешению: \emph{неизменяемость $\ne$ «нельзя удалять»}. Она означает
\emph{append-only, без тихих правок на месте}; само удаление — это \emph{тоже
appended, аудируемое событие}, а не нарушение дизайна. «Забыть» — первоклассная,
поддержанная операция, и реализуется тремя механизмами, которые работают вместе:

\begin{itemize}\itemsep2pt
\item \textbf{Граф (легко).} Убираем узел, ретрактим, перепропагируем. Каскад: удаление
исходного наблюдения понижает/сносит производные узлы, чьим \emph{провенансом} оно
было (ретракция по TMS, глава~\ref{ch:epistemics}). «Забудь X» — это ещё и «забудь
то, что из X выведено».
\item \textbf{Tombstone.} В лог дописывается delete-событие; при пересборке проекции
удалённое пропускается. Эффект немедленный, но сырое событие ещё \emph{в логе}
(годится для undo/аудита, но это не настоящее стирание).
\item \textbf{Crypto-shredding.} События (по сессии/скоупу) зашифрованы своим ключом;
«удалить» $=$ уничтожить ключ $\to$ шифртекст остаётся, но навсегда нечитаем.
\emph{Мгновенное, проверяемое} стирание без переписывания лога — идеально под «забудь
последние 10 минут» (эти минуты под эфемерным сессионным ключом; шреднул ключ —
стёр).
\item \textbf{Compaction.} Фоновая пересборка лога выбрасывает tombstoned-события —
сырое физически исчезает.
\end{itemize}

То есть «Рин, забудь X» $=$ tombstone $+$ каскадная ретракция производных $+$
шреддинг ключа (истинное стирание) $+$ компакция-уборка. Конкретно, а не «ну она как-
то забудет».

\begin{pitfall}[Crypto-shred рождает key management и требует taint-propagation]
Две дыры, без которых схема красива, но дырява. \emph{Первая}: шреддинг создаёт новый
центр тяжести — \textbf{управление ключами} (где ключи, кто их держит, гранулярность
сессия/скоуп). \emph{Вторая, важнее}: уничтожив ключ raw-события, ты \emph{не} стёр
след, если plaintext успел растечься — в debug-log, снапшот context-packer,
embedding-кэш, вектор-индекс, бэкап, telemetry-сводку или производный fusion-узел.
Поэтому рядом с «забудь» нужна дисциплина \textbf{taint-propagation}: всякий
производный объект наследует taint/провенанс от породившего события, а delete-cascade
чистит \emph{все} производные хранилища — граф, индексы, сводки, эмбеддинги, снапшоты,
pending-очереди, бэкапы. Иначе «забыл» будет только на бумаге.
\end{pitfall}

\begin{openq}[Гранулярность забвения]
На каком уровне держать ключи шифрования (событие / сессия / скоуп), что именно
каскадит при удалении (только прямые производные или транзитивно), и как показывать
пользователю, что забыто, — открыто. Но механика («забыть» = событие в логе $+$
shred $+$ каскад) зафиксирована.
\end{openq}

\section{Write-path: что попадает в память, а что нет}
\label{sec:writepath}

Отдельно проговорим, \emph{что именно} записывается, потому что «не писать» —
полноправный исход, а не сбой. Конвейер:
\begin{lstlisting}[style=pseudo]
raw_event -> parse/extract candidates
          -> [ epistemic_gate, salience_gate, anti_fusion_gate ]
          -> write | skip
          -> (if write) propagate
          -> audit_log   # и write, и skip логируются
\end{lstlisting}
\begin{itemize}\itemsep1pt
\item \textbf{epistemic-гейт}: стоит ли этому верить и с какой уверенностью (гл.~\ref{ch:epistemics}).
\item \textbf{salience-гейт}: значимо ли это для Рин (appraisal, гл.~\ref{ch:personality}).
\item \textbf{anti-fusion-гейт}: см.\ ниже.
\end{itemize}

\begin{principle}[response-useful association $\ne$ memory-worthy fusion]
Оператор должен уметь не только «нашёл красивую связь $\to$ породил узел», но и
явно \emph{не писать}: «со-активация есть, связного вывода нет», «связь эстетическая,
не эпистемическая», «это каламбур, не убеждение», «полезно для реплики, но не для
памяти». Рин может \emph{использовать} ассоциацию в ответе, не записывая её узлом —
иначе память зарастает «интересными мыслями», которые были хороши в моменте, но не
являются знанием.
\end{principle}

И \texttt{skip} обязательно логируется (хотя бы в debug): иначе не поймёшь, почему
Рин «не запомнила» важное.

\subsection{Коммит не бинарен: полосы, а не «только ночью»}
\label{sec:commitlanes}

Соблазнительно сказать: «walker'ы читают весь день, а пишет коммиттер только ночью».
Это \emph{слишком сильно} и ломает сознательное мышление. Если так сделать буквально,
Рин днём станет странной: ходит по графу, находит связи, что-то «понимает» — но не может
ни закрепить свежий внешний факт, ни воспользоваться дневной догадкой иначе как
эфемерным текстом. Правильный инвариант:

\begin{principle}[Коммит не бинарен]
Наблюдаемые события и рабочие гипотезы фиксируются \emph{немедленно}; долговременная
консолидация \emph{отложена}. Коммиттер один и сериализован, но работает не «ночью
против дня», а в нескольких \emph{полосах} по стоимости и durability.
\end{principle}

Полосы одного коммиттера:
\begin{itemize}\itemsep2pt
\item \textbf{immediate.} Пишется всегда и сразу: append в event-log (что произошло —
это \emph{журнал}, а не «убеждение»), прямые пользовательские наблюдения,
task-critical внешние факты, явные исправления, важные boundary\,/\,preference.
Откладывать это до ночи нельзя — потеряешь след происходящего.
\item \textbf{foreground pending.} Дневная рабочая память — это наш \emph{session-слой}
(гл.~\ref{ch:timescales}), не новая сущность: временные узлы-гипотезы, найденные
факты, цепочки рассуждения, кандидаты на консолидацию, \texttt{need\_to\_verify}. Живут
с TTL, низкой уверенностью, провенансом и статусом \texttt{pending} — уже объект,
которым можно \emph{пользоваться сейчас}, но ещё не вечная память.
\item \textbf{idle.} Микробатчи в простое: мелкий dedup, малые правки уверенности,
локальные проверки противоречий.
\item \textbf{sleep.} Тяжёлое: консолидация $A,B,C,D\to X$, массовое разрешение
конфликтов, fusion далёких ассоциаций, pruning, reindex, пересборка summary, ревизия
низкоуверенных потомков, promotion рекуррентных кандидатов.
\end{itemize}
Это и есть fast-path\,/\,slow-path: дешёвые $O(1)$-мутации (добавить узел-гипотезу,
инкремент) идут днём по быстрой полосе; дорогие $O(|V|\log|V|)$ и глобальные — ночью.
Критично: быстрая полоса — про \emph{латентность}, а не про поблажку. Она проходит тот
же write-path и те же гейты (epistemic\,/\,salience\,/\,anti-fusion,
раздел~\ref{sec:writepath}), просто синхронно.

Сводно, по трём процессам (и подчеркнём: строка записи \emph{не} ночецентрична — днём
работает immediate-lane):

\begin{center}
\begin{tabular}{p{0.19\textwidth} p{0.36\textwidth} p{0.36\textwidth}}
\toprule
\textbf{Процесс} & \textbf{Днём ($\rho\!\approx\!1$)} & \textbf{Фон\,/\,ночь ($\rho\!\approx\!0$)} \\
\midrule
Retrieval & локальный разреженный matvec-диффузия из HNSW-якорей (fast-path) & телепорт-скраб, эргодическое покрытие, большой $\alpha$ \\
Fusion\,/\,мышление & «мотыльки»: одиночные RW на остывающем $a_{\text{work}}$ $\to$ кандидаты в thought-inbox & честный рой (slow-path): треки, $n$-арные цепочки, далёкие мосты \\
Запись & \textbf{immediate-lane}: append в event-log, безопасные foreground-микрокоммиты (\texttt{observed}, \texttt{boundary}, правки \texttt{preference}); \texttt{pending} догадки — в thought-inbox & коммиттер (slow-path): судит гипотезы, тяжёлая консолидация, prune, reindex \\
\bottomrule
\end{tabular}
\end{center}

Так снимаются оба дневных противоречия. \emph{Сознательный поиск} не ждёт ночи: Рин
«гуглит здесь и сейчас», кладёт результат в event-log и session-граф как sourced
observation, использует в ответе — но со статусом «свежий источник, не консолидировано»,
а не «я давно это знаю». \emph{Догадки} тоже случаются днём — просто живут как активная
гипотеза, а взрослеют в durable-память позже.

\subsection{Thought-inbox: как дневная догадка попадает в реплику, не став фактом}
\label{sec:thoughtinbox}

Фоновый разведчик работает \emph{всегда} — даже в напряжённом диалоге, на остаточном
бюджете $B_{\text{bg}}=B_{\text{total}}(1-\rho)$ (гл.~\ref{ch:walker}), который никогда
не ноль. Поэтому «момент озарения» (связь House-Уилсон и Cast Away-Уилсон) может
родиться и днём. Но показывать Рин \emph{сырой след} walker'ов нельзя — это шум: сотни
посещённых узлов, случайные касания, мусорные мосты. Между следом и репликой стоит
фильтр — \emph{thought-inbox}: туда попадают не все касания, а только отобранные
предложения-инсайты (сильная, но неочевидная со-активация далёких узлов):
\begin{lstlisting}[style=pseudo]
type: insight_candidate
content: "Уилсон у House и Уилсон в Cast Away могут делить нарративную функцию"
status: unverified_hypothesis
salience: 0.72
provenance: [node_a, node_b, path_summary]
source: background_walker
ttl: 2h
recommended_action: mention_as_hypothesis | verify_later | ignore
\end{lstlisting}
Context-packer иногда подмешивает самые «горячие» из них в промпт — но \emph{с явной
эпистемической меткой} (это ровно наш packer-слой меток, раздел~\ref{sec:packer}):
«\emph{возможная мысль, не факт: между A и B может быть связь; статус — не проверено;
использовать как догадку, не как установленную память}». Дальше цикл замыкается сам:
если контекст диалога позволяет, Рин озвучивает догадку — и в \emph{этот} момент она
легализуется, падает в event-log и становится узлом со статусом
\texttt{need\_to\_verify}; если не пригодилась — тихо затухает по TTL и убирается
сборщиком. Изоляция ролей сохранена: walker \emph{читает и предлагает}, LLM
\emph{осмысляет и озвучивает}, коммиттер \emph{пишет}.

\begin{pitfall}[Не писать активацию на каждый шаг walker'а]
«Walker дёшев» — верно относительно LLM, но не бесплатно, и узкое место обычно не в
математике, а в обвязке (индексы, сериализация, локи, provenance-traversal,
scope-фильтры, гл.~\ref{ch:experiment}). Поэтому обходчик \emph{не} пишет активацию в
хранилище на каждом шаге: он копит visits\,/\,деньты активации локально в памяти
процесса и \emph{батчем} сбрасывает в transient-activation store (по-прежнему max, не
сумма). И на сбросе различает \texttt{access\_reason}: \texttt{user\_recall} укрепляет
след $B_{\text{base}}$, а \texttt{scrub\_visit} — нет.
\end{pitfall}

\section{Таксономия ретрива: не всё — блуждание}
\label{sec:retrievaltax}

Walker как «почти всё» — переобобщение. Не всякая операция памяти естественно
является ассоциативным блужданием:
\begin{itemize}\itemsep1pt
\item \textbf{associative} (якоря $+$ walkers) — для живости и ассоциаций (гл.~\ref{ch:walker}).
\item \textbf{forensic} — точный поиск по метаданным/провенансу/времени: «покажи всё
про $X$», «узлы с источником $Y$», «что противоречит этому факту», «последняя версия
prompt-модуля». Это индексы и query-planner, \emph{не} блуждание.
\item \textbf{creative} — walkers $+$ fusion $+$ сон (гл.~\ref{ch:sleep}).
\end{itemize}
Forensic-ретрив — не приятное дополнение, а \emph{необходимость} для обслуживания:
аудит, миграции, точное удаление конкретных узлов. Walker полноты не гарантирует, а
скучный индекс — гарантирует. Для живости Рин нужны walkers; для корректности памяти
нужны скучные индексы.

\section{Context-packer: как граф превращается в промпт}
\label{sec:packer}

Это интерфейс графа и LLM на инференсе — что именно \emph{видит} модель, и
\textbf{последний эпистемический гейт перед речью}. Если packer отдаёт модели плоский
список содержаний без статусов, вся дисциплина графа (уверенность, гипотезы,
депрецирование) теряется на последнем метре. Поэтому в контекст попадают \emph{не
только contents, но и метки}: \texttt{observed}, \texttt{hypothesis},
\texttt{deprecated}, \texttt{contested}, \texttt{low\_conf}, \texttt{time\_validity},
и при необходимости \texttt{do\_not\_state\_as\_fact}. Плюс провенанс для важных
утверждений и mood/persona-обусловливание. Тогда модель знает, что подать как факт,
что — как «похоже», а что не выдавать за истину вовсе. Держать компактно: короткий
контекст $=$ маленький KV-кэш $=$ меньше VRAM (раздел~\ref{sec:costs}).

Это \texttt{prompt context} — \emph{временная сборка} (граница из
главы~\ref{ch:architecture}), а не хранилище: спаковали, отдали модели, выбросили.

Здесь же — естественное место для \emph{eligibility-гейта}: перед упаковкой отсеять
узлы, которые \emph{есть в памяти, но не годятся для этого контекста} (например, не
класть в промпт то, что не стоит показывать при данном собеседнике или выносить в
публичный контент). Это шире удаления и \emph{у́же} тяжёлой governance-бюрократии:
не «кто владеет и на каких условиях», а простой guardrail
\emph{retrieval $\to$ eligibility-гейт $\to$ packer $\to$ output}. Тяжёлую систему
доступа/согласия мы сознательно не строим (вектор проекта сместился), но сам крючок
стоит держать в виду.

\begin{openq}[Eligibility как отложенный guardrail]
Полноценный слой видимости/раскрытия (что можно показать здесь, класть в промпт,
использовать в публичном контенте, обобщать) — отложен. Пока — минимальный гейт
«память может существовать, но быть негодной для данного контекста», без бюрократии.
Разворачивать ли его — решается, когда вернётся потребность в family/публичном режиме.
\end{openq}

\begin{openq}[Формат упаковки]
Как именно показывать модели статус узла (уверенность, гипотеза vs факт, провенанс),
сколько узлов включать и как их ранжировать под текущий запрос — открыто, меряется на
стенде.
\end{openq}

\section{Оценки стоимости: память, время, узкие места}
\label{sec:costs}

Оценки порядковые (для локального железа класса RTX 4090 / 8--24\,ГБ VRAM),
но они сразу показывают, где узкие места, а где нет.

\subsection{Стоимость операций}

\begin{center}
\begin{tabular}{p{0.30\textwidth} p{0.14\textwidth} p{0.46\textwidth}}
\toprule
\textbf{Операция} & \textbf{Порядок} & \textbf{Комментарий} \\
\midrule
LLM-инференс (композиция, объяснение) & \textbf{0,5--5 с} & \textbf{узкое место}; $\sim$2\,с/запрос на 7--8B локально \\
векторный поиск якорей (HNSW, $10^6$) & $\sim$мс & не узкое место \\
обход/spreading (сотни-тысячи узлов) & мс--10\,мс & C-ядро igraph, дёшево \\
пропагация уверенности (лавина) & $\sim$мс & дёшево \\
детекторы телеметрии & O(1) инкрем. & дёшево \\
распад/ACT-R & дёшево \emph{лениво} & полный проход по $N$ узлам — дорого; делать лениво (при доступе) или во сне, не на каждый тик \\
переваривание учебника & \textbf{минуты-часы} & одноразово; дорого из-за LLM-извлечения фактов; только в \textsc{Idle}/\textsc{Sleeping} \\
\bottomrule
\end{tabular}
\end{center}

\subsection{Латентность ответа}

Критический путь: \texttt{user\_msg} $\to$ поиск якорей (мс) $\to$ обход (мс) $\to$
композиция ответа (\textbf{LLM, $\sim$1--3\,с}) $\to$ ответ. \textbf{Бутылочное горло —
LLM-вызов}; всё графовое к латентности почти не добавляет. Цель — держать ответ в
$\sim$1--3\,с, то есть минимизировать \emph{число} LLM-вызовов на переднем плане (в
идеале один: граф уже подал контекст, модель только оформляет).

\subsection{Память}

\begin{itemize}\itemsep1pt
\item \textbf{VRAM} — модель (фиксировано, узкое место по размеру модели: 7--8B в
квантизации $\sim$5--8\,ГБ). Отсюда весь смысл «факты в граф, а не в веса».
\item \textbf{RAM} — векторный индекс горячего множества; для $10^6$ эмбеддингов
плоско $\sim$1,5\,ГБ, с HNSW/квантизацией меньше; \emph{тиринг} держит это ограниченным.
\item \textbf{Диск} — сам граф: $\sim$2--3\,КБ/узел, $10^5$--$10^6$ узлов $=$
0,3--3\,ГБ (глава~\ref{ch:experiment} для оценки роста). Дёшево; учебник $\approx 10^4$
узлов $\approx$ 30\,МБ.
\end{itemize}

\subsection{Где узкие места, а где нет}

\begin{principle}[Всё сводится к LLM-слотам]
Дорог только GPU/LLM — и по латентности ответа, и по фоновой пропускной способности.
Всё графовое (хранилище, обход, статистика, распад, векторный поиск) — дёшево (мс) и
узким местом \emph{не} является. Поэтому вся инженерная задача рантайма сводится к
одному: \emph{как распределить дефицитные LLM-вызовы} между быстрым ответом (передний
план) и фоном (fusion/метакогниция/переваривание). Это и делает планировщик
(раздел~\ref{sec:scheduler}).
\end{principle}

\begin{openq}[Цели по латентности и throughput]
Конкретные цели (ответ $<$ N\,с; сколько фоновых fusion в минуту в простое; когда
переключать модель на меньшую ради скорости) — под стенд и под конкретное железо.
Динамический роутинг «мелкая модель на рутину, крупная на метакогницию» (как у Toy и
др.) — очевидный кандидат оптимизации.
\end{openq}
